\section{Superposition}

\begin{exercise}[The monochord]
\textbf{Experiment.}
A string fixed at both ends rings only at frequencies
\[
        f_n=nf_1,\qquad n=1,2,\ldots,
\]
and an arbitrary small displacement separates into the corresponding normal
modes.  Let \(V\) be the real vector space of displacements, let
\(K:V\to V^*\) encode the restoring force, and let
\(M:V\to V^*\) encode kinetic energy.  Assume that \(M\) is positive and that
\(K\) is symmetric with respect to the natural pairing.

Explain the observed decomposition without choosing coordinates.  Show that a
mode of angular frequency \(\omega\) is a nonzero solution of
\[
        Kv=\omega^2Mv.
\]
Use the \(M\)-inner product to explain why modes with distinct frequencies are
orthogonal.  For an ideal string, identify the boundary condition that gives
\(\omega_n=n\omega_1\), and state which assumptions fail when the amplitude is
large enough that superposition is no longer observed.
\end{exercise}

\begin{exercise}[Young's fringes]
\textbf{Experiment.}
Young's two-slit experiment produces alternating bright and dark bands
\cite{young}.  Blocking either slit removes the bands.  At a detector point,
let the two alternatives contribute \(a,b\in\C\), while the measured intensity
is proportional to \(|a+b|^2\).

Explain why adding measured intensities cannot produce the dark bands.
Expand \(|a+b|^2\), isolate the interference term, and show that it depends on
the relative phase but not on a simultaneous multiplication
\((a,b)\mapsto(e^{i\theta}a,e^{i\theta}b)\).  Reformulate the physical state as
a ray in a complex vector space.  Explain why the quotient by nonzero scalar
multiplication records the experimental invariance more faithfully than any
choice of phase convention.
\end{exercise}

\begin{exercise}[One photon, two paths]
\textbf{Experiment.}
In a balanced Mach--Zehnder interferometer, single photons arrive one at a
time, yet one output can be dark after many trials.  Inserting a phase
\(\phi\) in one arm changes the output probabilities to
\[
        p_0=\cos^2(\phi/2),\qquad p_1=\sin^2(\phi/2).
\]
Let \(H\) be the two-dimensional path space and let a lossless beam splitter
be represented by a unitary map \(B:H\to H\).

Find a coordinate-free expression for the phase shifter using the two path
subspaces and their orthogonal projections.  Compose two beam splitters with
this map and recover the observed probabilities.  Explain why a classical
mixture of \emph{upper path} and \emph{lower path} cannot make one output
dark, and why the experiment nevertheless does not determine a preferred
basis for \(H\).
\end{exercise}

\begin{exercise}[Stern--Gerlach splitting]
\textbf{Experiment.}
A beam of silver atoms passing through an inhomogeneous magnetic field splits
into two beams, not a continuous smear \cite{stern-gerlach}.  Repeating the
experiment with a rotated magnet changes the relative counts.  Model a
measurement direction \(n\in\R^3\) by a self-adjoint operator \(S_n\) on a
two-dimensional complex Hilbert space.

Find structural assumptions on \(n\mapsto S_n\) which force
\(S_n^2=(\hbar^2/4)\id\).  Deduce the two observed values
\(\pm\hbar/2\).  If \(P_n^+\) and \(P_m^+\) are the positive-beam projections,
show that the probability of passing both filters is
\(\Tr(P_n^+P_m^+)=(1+n\cdot m)/2\).  Explain the disappearance and return of a
blocked beam when an intermediate, differently oriented analyzer is inserted.
\end{exercise}

\begin{exercise}[Tomography of an open evolution]
\textbf{Experiment.}
Quantum process tomography reconstructs laboratory evolutions which are not
unitary: attenuation, dephasing, and relaxation all occur
\cite{altepeter}.  Let \(\Phi:\operatorname{End}(H)\to\operatorname{End}(K)\)
be the reconstructed linear map.  Experiments require density operators to
remain positive even when the system is entangled with an untouched auxiliary
system \(A\).

Show that the requirement is positivity of
\(\Phi\otimes\id_{\operatorname{End}(A)}\) for every finite-dimensional
\(A\).  Prove that it suffices to test \(A\cong H\).  From the positive
operator obtained by applying \(\Phi\otimes\id\) to a maximally entangled
state, derive operators \(V_i:H\to K\) satisfying
\[
        \Phi(\rho)=\sum_i V_i\rho V_i^*.
\]
Explain why preservation of total detection probability is equivalent to
\(\sum_iV_i^*V_i=\id_H\).
\end{exercise}

\begin{exercise}[Chladni figures]
\textbf{Experiment.}
Sand on a driven metal plate collects along sharp nodal curves, and the
pattern changes abruptly as the driving frequency is varied.
Let \(V\) be the displacement space and let the plate's elastic and inertial
forms be \(K\) and \(M\).

Explain why persistent patterns occur near the spectrum of the generalized
operator \(K-\omega^2M\).  Show that the sand records the zero set of an
eigenvector rather than its coordinates.  Explain why a geometrically
symmetric plate can have degenerate modes, why a small clamp or added mass
splits them, and why an arbitrary linear combination is observed when the
driving frequency meets an unsplit eigenspace.
\end{exercise}

\begin{exercise}[Malus's law]
\textbf{Experiment.}
Light transmitted through two ideal linear polarizers has intensity
\(I=I_0\cos^2\theta\), where \(\theta\) is their relative angle.
Represent a pure polarization by a ray in a two-dimensional complex inner
product space and a polarizer by an orthogonal projection \(P\) of rank one.

Explain the observed law by showing that the transmitted fraction from ray
\([v]\) is \(\|Pv\|^2/\|v\|^2\).  Prove that this number depends only on the
two rays and not on phases or chosen bases.  Explain why inserting a third
polarizer between crossed polarizers can restore nonzero transmission.
\end{exercise}

\begin{exercise}[Ramsey fringes]
\textbf{Experiment.}
Two separated resonant pulses applied to an atom produce an oscillating
excitation probability as the pulse separation or oscillator frequency is
varied.  Let the two energy eigenspaces be orthogonal lines in \(H\); model
each pulse by a unitary mixing them and free evolution by the spectral
projections of the Hamiltonian.

Compose the three maps and recover a sinusoidal detection probability.
Explain why the phase accumulated during the unobserved interval affects the
final population.  Identify the invariant relative phase and explain how the
fringe period measures the energy splitting without selecting coordinates on
\(H\).
\end{exercise}

\begin{exercise}[Neutron diffraction]
\textbf{Experiment.}
A monochromatic neutron beam scattered by a crystal gives intense peaks only
at particular angles.  Associate an amplitude to scattering from each lattice
site and let translations act on the wave's one-dimensional phase space.

Explain why the total amplitude is a sum of translation characters.  Show
that constructive interference occurs precisely when the momentum transfer
defines the trivial character on the crystal lattice.  Recover Bragg's law
after choosing a family of lattice planes.  Explain why the reciprocal
lattice, defined as a character group, is independent of any chosen primitive
cell.
\end{exercise}

\begin{exercise}[Interference of large molecules]
\textbf{Experiment.}
Molecular beams passing through a diffraction grating produce interference
fringes even when the molecules have many internal degrees of freedom.
Associate a complex amplitude to each center-of-mass path and an internal
state to the molecule.

Show that the detected intensity contains cross terms weighted by the inner
product of the internal states carried along the two paths.  Explain why
unobserved internal excitations can reduce visibility without selecting a
definite trajectory.  Identify the coherence condition which lets a composite
object display the same superposition law as a photon.
\end{exercise}
